% =============================================================================
%  圆盘质量不平衡引发的脉动转矩：MathType 可直接读取的 LaTeX 片段
%  版本：2026-11-13
%
%  用法（MathType）：
%    1) MathType → Preferences → Workspace Preferences
%       勾选 "Allow TeX language input"（允许 TeX 语言输入）；
%    2) 复制下面任意一行的公式（不含行首的 % 注释行）；
%    3) 在 MathType 编辑窗口中粘贴并回车，即转成公式。
%
%  约定：每条公式单行、独立、可单独粘贴；不使用环境、编号标签、
%        \text 类命令与自定义宏。
%  核心结论：水平轴不平衡的绕轴脉动转矩来自重力，幅值 M_t g e（与转速无关）；
%            m e ω^2 是径向离心力，不是绕轴转矩。
% =============================================================================

% U1 总质量
M_t=M+m

% U2 质心偏心距
e=\frac{mR}{M_t}

% U3 绕轴转动惯量
I_n=\frac{1}{2}MR^2+I_s+mR^2

% U4 多质量等效偏心
e=\frac{1}{M_t}\left|\sum_i m_i\vec{r}_i\right|

% U5 多质量转动惯量
I_n=\frac{1}{2}MR^2+\sum_i m_iR^2

% U6 转角
\theta=\omega t+\theta_0

% U7 质心位置
\vec{r}_C=e\left(\sin\theta,\ -\cos\theta\right)

% U8 重力矩
\tau_g=\left(\vec{r}_C\times\vec{F}_g\right)_Z=-M_t g e\sin\theta

% U9 用边缘质量表示
\tau_g=-m g R\sin\theta

% U10 重力矩幅值
T_g=M_t g e=m g R

% U11 时域脉动转矩
\tau_g\left(t\right)=-T_g\sin\left(\omega t+\theta_0\right)

% U12 轴倾角修正
T_g=M_t g e\cos\gamma

% U13 离心力
F_u=M_t e\omega^2=mR\omega^2

% U14 离心力对轴无力矩
\tau_Z\left(\vec{F}_u\right)=\hat{Z}\cdot\left(\vec{r}_C\times\vec{F}_u\right)=0

% U15 轴承 A 反力
\vec{R}_A=\frac{b}{L}\left(\vec{F}_u+M_t\vec{g}\right)

% U16 轴承 B 反力
\vec{R}_B=\frac{a}{L}\left(\vec{F}_u+M_t\vec{g}\right)

% U17 绕轴转矩方程
I_n\dot{\omega}=T_e-T_L\left(\theta\right)-T_f\left(\theta,\omega\right)+\tau_g\left(\theta\right)

% U18 稳态电磁转矩
T_e\approx T_L+T_f+M_t g e\sin\theta

% U19 转速纹波
\Delta\omega_1=\left|S\left(j\omega\right)\right|\frac{M_t g e}{I_n\omega}

% U20 法向载荷平方
\left|\vec{F}_n\right|^2=\left(M_t g\right)^2+\left(mR\omega^2\right)^2+2M_t g\,mR\omega^2\cos\varphi

% U21 摩擦纹波
\Delta T_f=\mu_b r_b\frac{M_t g\,mR\omega^2}{\sqrt{\left(M_tg\right)^2+\left(mR\omega^2\right)^2}}

% U22 动态不平衡力偶
M_u=mR\omega^2 l

% U23 力偶的轴承载荷
R_{uA}=R_{uB}=\frac{M_u}{L}

% U24 转矩与力的换算
T_g=\frac{g}{\omega^2}F_u

% U25 反解离心力
F_u=\frac{\omega^2}{g}T_g

% U26 由型线反推
F_u=\frac{T_1\omega^2}{g}

% U27 台架反推数值
F_u=\frac{0.0571\times\left(62.83\right)^2}{9.81}\approx23\ \mathrm{N}

% U28 等效不平衡量
mR=\frac{F_u}{\omega^2}\approx5.8\times10^{-3}\ \mathrm{kg\cdot m}

% U29 现实不平衡量下的离心力
F_u=mR\omega^2\approx0.39\ \mathrm{N}

% U30 相应重力矩
T_g=\frac{g}{\omega^2}F_u\approx9.8\times10^{-4}\ \mathrm{N\cdot m}

% U31 占比
\frac{T_g}{T_1}\approx1.7\%

% U32 竖直轴重力矩
T_g=0

% U33 质心轴向偏置的附加力偶
\pm\frac{z_C\hat{Z}\times\vec{F}_u}{L}

